\chapter{Скрученная область определения полного тензорного гессиана}
\label{ch:v6-bosonic-defect-full-tensor-polar-hessian-gate}

\section{Почему нельзя просто добавить компоненты}

Предыдущий гейт закрыл все целые угловые числа эффективной
пятикомпонентной подсистемы. Для полного набора $Q+T+B$ буквальное повторение
того же ряда Фурье неверно. Фон вихря имеет голономию
\begin{equation}
 g=\exp(2\pi J_0/3)\in Z_3,
\end{equation}
поэтому сопутствующий внутренний базис после обхода нити возвращается с
действием $g$, а не тождественно.

Это уточняет статус старого тринадцатикомпонентного радиального скана. Он
остаётся локальной проверкой робастности потенциального и радиального блоков,
но не является спектром глобальных однозначных сечений вокруг нити. Требование
учитывать монодромию поля в присутствии нетривиальной голономии является
стандартным для вихревых и поверхностных дефектов
\cite{full-tensor-bucher-alice,full-tensor-gomis-surface}.

\section{Весовое разложение}

Положим $h_0=J_0/3$. Целые внутренние веса оператора $3(-i\rho(h_0))$
равны
\begin{align}
 V_2 &: (-2,-1,0,1,2),\\
 V_3 &: (-3,-2,-1,0,1,2,3),\\
 \mathfrak{so}(3)_{\mathrm{ad}} &: (-1,0,1).
\end{align}
Остаточный элемент различает веса только по модулю три. В каждом характере
$c=-1,0,1$ находятся четыре комплексные компоненты материи, одна радиальная
и одна угловая компонента связности. Поэтому полный полярный гессиан
распадается на три шестикомпонентных комплексных радиальных семейства.

Пусть $s_g\in\{-1,0,1\}$ --- вес компоненты связности выбранного характера,
а $j\in\mathbb Z$. Общая угловая экспонента в сопутствующем базисе и
однозначная лабораторная гармоника компоненты веса $s$ равны
\begin{equation}
 \mu=j-\frac{s_g}{3},
 \qquad
 n_s=\mu+\frac{s}{3}
     =j+\frac{s-s_g}{3}\in\mathbb Z.
 \label{eq:v6-full-tensor-twisted-harmonic-rule}
\end{equation}
Именно $n_s$, а не общий формальный $m$, задаёт регулярность каждой
компоненты в сердцевине.

\section{Машинная проверка замкнутости секторов}

Проверка выполнена на полном двенадцатимерном потенциальном гессиане
пяти компонент $Q$ и семи компонент $T$. На четырёх точках вдоль профиля
максимальный остаток ковариантности относительно $Z_3$ равен
\begin{equation}
 2.17\cdot10^{-6},
\end{equation}
что соответствует погрешности конечного дифференцирования потенциала.
Для отображения из трёх компонент связности в касательные к калибровочной
орбите остаток равен
\begin{equation}
 3.97\cdot10^{-15}.
\end{equation}
Следовательно, потенциальный блок и сцепление связности действительно
сохраняют три остаточных характера и могут собираться независимо.

Если оставить только общий целочисленный ряд в сопутствующем базисе, глобально
допустимым будет лишь нейтральный характер: шесть комплексных радиальных
полей из восемнадцати. Двенадцать полей двух скрученных характеров будут
потеряны. Это не малая поправка, а две трети полного оператора.

Переносные моды лежат в нейтральном характере при $j=\pm1$; после сборки
спектра там должны возникнуть ровно два вещественных нуля.

\section{Нулевая гипотеза и стоп-критерий}

Нулевая гипотеза следующего гейта: хотя эффективный нейтральный блок
положителен, один из двух скрученных тензорных секторов содержит отрицательную
моду. Для проверки требуется собрать квадратичную форму с
$Z_Q=Z_T=1/3$, всеми тремя компонентами связности, фоновой калибровкой и
условиями регулярности по правилу~\eqref{eq:v6-full-tensor-twisted-harmonic-rule}.

Стоп-критерий: сходящаяся отрицательная мода любого характера закрывает
прямую вихревую нить. Если же численный оператор использует общую
целочисленную гармонику для компонент разных внутренних весов, такой расчёт
не считается проверкой полного гессиана.

\begin{protocol}
\begin{enumerate}
 \item полное число комплексных радиальных полей: \textbf{18};
 \item число остаточных характеров: \textbf{три};
 \item размер каждого радиального семейства: \textbf{шесть};
 \item обычный общий целочисленный ряд: \textbf{отвергнут};
 \item потеря при нескрученном ограничении: \textbf{12 из 18 полей};
 \item глобальная область определения полного гессиана: \textbf{выведена};
 \item полный тензорный спектр: \textbf{не вычислен};
 \item полная устойчивость нити: \textbf{не закрыта};
 \item следующий гейт: \textbf{спектр трёх скрученных секторов}.
\end{enumerate}
\end{protocol}

Машинный сертификат сохранён в
\path{s2t_v6_bosonic_defect_full_tensor_polar_hessian_gate_results.json}.

\begin{thebibliography}{9}
\bibitem{full-tensor-bucher-alice}
M.~Bucher, H.-K.~Lo, J.~Preskill,
\emph{Topological approach to Alice electrodynamics},
arXiv:hep-th/9112039.

\bibitem{full-tensor-gomis-surface}
J.~Gomis, S.~Matsuura,
\emph{Holography of BPS surface operators},
arXiv:0812.1420.
\end{thebibliography}